\section{Determinants and moment-ratio tools}
Moment-ratio and determinant inequalities provide a compressed way to propagate positivity through derivative hierarchies.  These tools are particularly useful for Mills-ratio and Bessel-type estimates where the same normal form controls many levels at once.
\begin{theorem}[Mills ratio positive Laplace representation and complete monotonicity]
\label{res:from-mills-laplace-cm-normal-form}
For \(t>0\), the standard normal Mills ratio satisfies \(r(t)=\int_0^\infty e^{-tu-u^2/2}\,du\). Hence \(r\) is completely monotone on \((0,\infty)\).
\end{theorem}
\begin{proof}
For the standard normal Mills ratio, the normalizing constants cancel:
\[
r(t)=\frac{\int_t^\infty e^{-z^2/2}\,dz}{e^{-t^2/2}}
=e^{t^2/2}\int_t^\infty e^{-z^2/2}\,dz.
\]
Put \(z=t+u\). Then
\[
r(t)=\int_0^\infty e^{-tu-u^2/2}\,du.
\]
For \(t>0\), differentiating under the integral gives
\[
(-1)^n r^{(n)}(t)=\int_0^\infty u^n e^{-tu-u^2/2}\,du>0.
\]
Thus \(r\) is completely monotone on \((0,\infty)\).

From the same integral,
\[
r'(t)=-\int_0^\infty u e^{-tu-u^2/2}\,du.
\]
Since
\[
\frac{d}{du}e^{-tu-u^2/2}=-(t+u)e^{-tu-u^2/2},
\]
integration over \([0,\infty)\) gives
\[
\int_0^\infty (t+u)e^{-tu-u^2/2}\,du=1.
\]
Therefore \(\int_0^\infty u e^{-tu-u^2/2}\,du=1-tr(t)\), and
\[
r'(t)=tr(t)-1.
\]

Set \(P_0=1\), \(Q_0=0\), and suppose
\[
r^{(n)}(t)=P_n(t)r(t)+Q_n(t).
\]
Differentiating and using \(r'=tr-1\),
\[
r^{(n+1)}(t)=(P_n'(t)+tP_n(t))r(t)+(Q_n'(t)-P_n(t)).
\]
Thus
\[
P_{n+1}=P_n'+tP_n,\qquad Q_{n+1}=Q_n'-P_n.
\]
This inductively gives the claimed polynomial recurrence.

For each integer \(L\ge0\), define
\[
f_L(t)=(-1)^L r^{(L)}(t)=\int_0^\infty u^L e^{-tu-u^2/2}\,du.
\]
This is a positive Laplace transform. Let
\[
a_j(t)=(-1)^j f_L^{(j)}(t)=\int_0^\infty u^{L+j}e^{-tu-u^2/2}\,du.
\]
Normalize the positive measure \(u^L e^{-tu-u^2/2}du\) to a probability law \(X\). Up to the positive factor \(a_0(t)^2\), the determinant expression is
\[
\mathbb E[X^4]-4\mathbb E[X^3]\mathbb E[X]+3\mathbb E[X^2]^2.
\]
With \(\mu=\mathbb E[X]\), this equals
\[
\mathbb E[(X-\mu)^4]+3\operatorname{Var}(X)^2\ge0.
\]
Hence
\[
f_L^{(4)}f_L-4f_L^{(3)}f_L'+3(f_L'')^2\ge0.
\]
Since \(f_L^{(j)}=(-1)^L r^{(L+j)}\), this is exactly
\[
r^{(L+4)}r^{(L)}-4r^{(L+3)}r^{(L+1)}+3(r^{(L+2)})^2\ge0.
\]
\end{proof}

\begin{theorem}[Yang Tian Bessel W general zero partial fraction proves W\_nu sqrt is Bernstein]
\label{res:yt-bessel-w-general-zero-partial-fraction}
For every \(\nu>-1\), if \(j_{\nu+1,n}\) are the positive zeros of \(J_{\nu+1}\), then \(G_\nu(s)=W_\nu(\sqrt{s})\) satisfies \(G_\nu(s)=2(\nu+1)+2\sum_{n\ge1}s/(s+j_{\nu+1,n}^2)\) and \(G_\nu'(s)=2\sum_{n\ge1}j_{\nu+1,n}^2/(s+j_{\nu+1,n}^2)^2\). Consequently \(G_\nu\) is a Bernstein function on \((0,\infty)\).
\end{theorem}
\begin{proof}
Put \(\mu=\nu+1>0\). The canonical product for the modified Bessel function gives, locally uniformly in \(z\),
\[
I_\mu(z)=\frac{(z/2)^\mu}{\Gamma(\mu+1)}
\prod_{n=1}^{\infty}\left(1+\frac{z^2}{j_{\mu,n}^2}\right),
\]
where \(j_{\mu,n}\) are the positive zeros of \(J_\mu\). Therefore
\[
z\frac{I_\mu'(z)}{I_\mu(z)}
=\mu+2\sum_{n=1}^{\infty}\frac{z^2}{z^2+j_{\mu,n}^2}.
\]
The recurrence
\[
I_\mu'(z)=I_{\mu-1}(z)-\frac{\mu}{z}I_\mu(z)
\]
then yields
\[
W_\nu(z)
=z\frac{I_{\mu-1}(z)}{I_\mu(z)}
=2\mu+2\sum_{n=1}^{\infty}\frac{z^2}{z^2+j_{\mu,n}^2}.
\]

Set \(s=z^2\). Then
\[
G_\nu(s):=W_\nu(\sqrt s)
=2(\nu+1)+2\sum_{n=1}^{\infty}\frac{s}{s+j_{\nu+1,n}^2}.
\]

Termwise differentiation is justified because \(j_{\nu+1,n}\sim \pi n\), so the differentiated sums converge locally uniformly on \((0,\infty)\). Thus
\[
G_\nu'(s)
=2\sum_{n=1}^{\infty}
\frac{j_{\nu+1,n}^2}{(s+j_{\nu+1,n}^2)^2}.
\]
For every \(k\ge0\),
\[
(-1)^k\frac{d^k}{ds^k}
\frac{j_{\nu+1,n}^2}{(s+j_{\nu+1,n}^2)^2}
=
(k+1)!\frac{j_{\nu+1,n}^2}{(s+j_{\nu+1,n}^2)^{k+2}}
\ge0.
\]
The locally uniformly convergent sum preserves these inequalities, hence \(G_\nu'\) is completely monotone and \(G_\nu\) is a Bernstein function.

If \(0<\tau\le1/2\), then \(x^{2\tau}\) is a Bernstein function. By Bernstein-function composition closure,
\[
x\mapsto W_\nu(x^\tau)=G_\nu(x^{2\tau})
\]
is a Bernstein function on \((0,\infty)\).
\end{proof}

\begin{proposition}[positive Laplace tilted moments imply strict monotonicity of a\_{n+1}/(a\_n a\_{n+2})]
\label{res:cm-laplace-moment-ratio-monotonicity}
Let \(a_j(x)=\int_0^\infty t^j e^{-xt}\,d\mu(t)>0\) be tilted moments of a positive Laplace measure for the needed indices. Then, for every \(n\ge0\), \(B_n(x)=a_{n+1}(x)/(a_n(x)a_{n+2}(x))\) is strictly increasing on its domain.
\end{proposition}
\begin{proof}
Let \(\mu\) be a positive Borel measure on \([0,\infty)\) with finite moments after the tilt \(e^{-xt}\). Define
\[
a_j(x)=\int_0^\infty t^j e^{-xt}\,d\mu(t)>0
\]
for the needed indices. For fixed \(n\ge0\), set
\[
B_n(x)=\frac{a_{n+1}(x)}{a_n(x)a_{n+2}(x)}.
\]
Then \(B_n\) is strictly increasing whenever \(a_{n+1}(x)>0\).

Indeed \(a_j'(x)=-a_{j+1}(x)\). Put
\[
r_j(x)=\frac{a_{j+1}(x)}{a_j(x)}.
\]
Moment log-convexity, equivalently Cauchy--Schwarz applied to
\[
\int t^j e^{-xt}\,d\mu(t),
\]
gives \(r_{j+1}(x)\ge r_j(x)\). Differentiating,
\[
\frac{d}{dx}\log B_n(x)
=-\frac{a_{n+2}}{a_{n+1}}
+\frac{a_{n+1}}{a_n}
+\frac{a_{n+3}}{a_{n+2}}
=r_n(x)-r_{n+1}(x)+r_{n+2}(x).
\]
Since \(r_{n+2}(x)\ge r_{n+1}(x)\) and \(r_n(x)>0\), the logarithmic derivative is positive. Hence \(B_n\) is strictly increasing.

Let
\[
f_k(x)=x\beta_k(x).
\]
From the integral representation of \(\beta_k\), integration by parts gives
\[
f_k(x)
=x\int_0^\infty e^{-xt}\frac{dt}{1+e^{-kt}}
=\frac12+\int_0^\infty e^{-xt}\frac{k e^{-kt}}{(1+e^{-kt})^2}\,dt.
\]
Thus \(f_k\) is the Laplace transform of a positive measure consisting of an atom \(1/2\) at \(0\) plus the positive density
\[
w_k(t)=\frac{k e^{-kt}}{(1+e^{-kt})^2}
\]
on \((0,\infty)\). All tilted moments are finite, and \(a_{j}(x)>0\) for \(j\ge1\).

For \(a_j(x)=(-1)^j f_k^{(j)}(x)\), the bridge lemma gives that
\[
B_{k,n}(x)=\frac{a_{n+1}(x)}{a_n(x)a_{n+2}(x)}
\]
is strictly increasing for every \(k>0\), \(n\ge0\), and \(x>0\). Since
\[
R_{k,n}(x)=(-1)^{n+1}B_{k,n}(x),
\]
we get the precise monotonicity law:

if \(n\) is odd, \(R_{k,n}\) is strictly increasing on \((0,\infty)\);
if \(n\) is even, \(R_{k,n}\) is strictly decreasing on \((0,\infty)\).

This solves the Yin--Zhang Nielsen \(k\)-beta derivative-ratio open problem in the parity-refined form.
\end{proof}

\begin{theorem}[Mills ratio all-L fourth-order determinant inequality]
\label{res:from-mills-derivative-determinant-all-l}
For every integer \(L\ge0\) and every \(t>0\), \(r^{(L+4)}(t)r^{(L)}(t)-4r^{(L+3)}(t)r^{(L+1)}(t)+3(r^{(L+2)}(t))^2\ge0\), where \(r\) is the standard normal Mills ratio.
\end{theorem}
\begin{proof}
For the standard normal Mills ratio, the normalizing constants cancel:
\[
r(t)=\frac{\int_t^\infty e^{-z^2/2}\,dz}{e^{-t^2/2}}
=e^{t^2/2}\int_t^\infty e^{-z^2/2}\,dz.
\]
Put \(z=t+u\). Then
\[
r(t)=\int_0^\infty e^{-tu-u^2/2}\,du.
\]
For \(t>0\), differentiating under the integral gives
\[
(-1)^n r^{(n)}(t)=\int_0^\infty u^n e^{-tu-u^2/2}\,du>0.
\]
Thus \(r\) is completely monotone on \((0,\infty)\).

From the same integral,
\[
r'(t)=-\int_0^\infty u e^{-tu-u^2/2}\,du.
\]
Since
\[
\frac{d}{du}e^{-tu-u^2/2}=-(t+u)e^{-tu-u^2/2},
\]
integration over \([0,\infty)\) gives
\[
\int_0^\infty (t+u)e^{-tu-u^2/2}\,du=1.
\]
Therefore \(\int_0^\infty u e^{-tu-u^2/2}\,du=1-tr(t)\), and
\[
r'(t)=tr(t)-1.
\]

Set \(P_0=1\), \(Q_0=0\), and suppose
\[
r^{(n)}(t)=P_n(t)r(t)+Q_n(t).
\]
Differentiating and using \(r'=tr-1\),
\[
r^{(n+1)}(t)=(P_n'(t)+tP_n(t))r(t)+(Q_n'(t)-P_n(t)).
\]
Thus
\[
P_{n+1}=P_n'+tP_n,\qquad Q_{n+1}=Q_n'-P_n.
\]
This inductively gives the claimed polynomial recurrence.

For each integer \(L\ge0\), define
\[
f_L(t)=(-1)^L r^{(L)}(t)=\int_0^\infty u^L e^{-tu-u^2/2}\,du.
\]
This is a positive Laplace transform. Let
\[
a_j(t)=(-1)^j f_L^{(j)}(t)=\int_0^\infty u^{L+j}e^{-tu-u^2/2}\,du.
\]
Normalize the positive measure \(u^L e^{-tu-u^2/2}du\) to a probability law \(X\). Up to the positive factor \(a_0(t)^2\), the determinant expression is
\[
\mathbb E[X^4]-4\mathbb E[X^3]\mathbb E[X]+3\mathbb E[X^2]^2.
\]
With \(\mu=\mathbb E[X]\), this equals
\[
\mathbb E[(X-\mu)^4]+3\operatorname{Var}(X)^2\ge0.
\]
Hence
\[
f_L^{(4)}f_L-4f_L^{(3)}f_L'+3(f_L'')^2\ge0.
\]
Since \(f_L^{(j)}=(-1)^L r^{(L+j)}\), this is exactly
\[
r^{(L+4)}r^{(L)}-4r^{(L+3)}r^{(L+1)}+3(r^{(L+2)})^2\ge0.
\]
\end{proof}

The next result applies the preceding method to the fixed source problem \textbf{APP-0040}.
\begin{corollary}[APP-0040: Yang-Tian Bessel-\(W\) power-Bernstein conjecture is true]
\label{res:yt-bessel-w-power-bernstein-conjecture}
For \(\tau\in(0,1/2]\) and \(\nu>-1\), prove that \(x\mapsto W_\nu(x^\tau)\), where \(W_\nu(x)=xI_\nu(x)/I_{\nu+1}(x)\), is a Bernstein function on \((0,\infty)\).

This resolves the fixed application label \textbf{APP-0040}: Solve Yang and Tian's Bessel-\(W\) power-Bernstein conjecture on the source interval.
\emph{Source reference.} \cite{app-app-0040-source}.
\end{corollary}
\begin{proof}
Put \(\mu=\nu+1>0\). The canonical product for the modified Bessel function gives, locally uniformly in \(z\),
\[
I_\mu(z)=\frac{(z/2)^\mu}{\Gamma(\mu+1)}
\prod_{n=1}^{\infty}\left(1+\frac{z^2}{j_{\mu,n}^2}\right),
\]
where \(j_{\mu,n}\) are the positive zeros of \(J_\mu\). Therefore
\[
z\frac{I_\mu'(z)}{I_\mu(z)}
=\mu+2\sum_{n=1}^{\infty}\frac{z^2}{z^2+j_{\mu,n}^2}.
\]
The recurrence
\[
I_\mu'(z)=I_{\mu-1}(z)-\frac{\mu}{z}I_\mu(z)
\]
then yields
\[
W_\nu(z)
=z\frac{I_{\mu-1}(z)}{I_\mu(z)}
=2\mu+2\sum_{n=1}^{\infty}\frac{z^2}{z^2+j_{\mu,n}^2}.
\]

Set \(s=z^2\). Then
\[
G_\nu(s):=W_\nu(\sqrt s)
=2(\nu+1)+2\sum_{n=1}^{\infty}\frac{s}{s+j_{\nu+1,n}^2}.
\]

Termwise differentiation is justified because \(j_{\nu+1,n}\sim \pi n\), so the differentiated sums converge locally uniformly on \((0,\infty)\). Thus
\[
G_\nu'(s)
=2\sum_{n=1}^{\infty}
\frac{j_{\nu+1,n}^2}{(s+j_{\nu+1,n}^2)^2}.
\]
For every \(k\ge0\),
\[
(-1)^k\frac{d^k}{ds^k}
\frac{j_{\nu+1,n}^2}{(s+j_{\nu+1,n}^2)^2}
=
(k+1)!\frac{j_{\nu+1,n}^2}{(s+j_{\nu+1,n}^2)^{k+2}}
\ge0.
\]
The locally uniformly convergent sum preserves these inequalities, hence \(G_\nu'\) is completely monotone and \(G_\nu\) is a Bernstein function.

If \(0<\tau\le1/2\), then \(x^{2\tau}\) is a Bernstein function. By Bernstein-function composition closure,
\[
x\mapsto W_\nu(x^\tau)=G_\nu(x^{2\tau})
\]
is a Bernstein function on \((0,\infty)\).
\end{proof}

\begin{lemma}[Yang Tian Bessel W\_nu x power tau Bernstein conjecture for tau <= 1/2 and nu > -1]
\label{res:yang-tian-bessel-w-bernstein-conjecture}
For \(\tau\in(0,1/2]\) and \(\nu>-1\), the function \(x\mapsto W_\nu(x^\tau)\), where \(W_\nu(x)=xI_\nu(x)/I_{\nu+1}(x)\), is a Bernstein function on \((0,\infty)\).
\end{lemma}
\begin{proof}
Put \(\mu=\nu+1>0\). The canonical product for the modified Bessel function gives, locally uniformly in \(z\),
\[
I_\mu(z)=\frac{(z/2)^\mu}{\Gamma(\mu+1)}
\prod_{n=1}^{\infty}\left(1+\frac{z^2}{j_{\mu,n}^2}\right),
\]
where \(j_{\mu,n}\) are the positive zeros of \(J_\mu\). Therefore
\[
z\frac{I_\mu'(z)}{I_\mu(z)}
=\mu+2\sum_{n=1}^{\infty}\frac{z^2}{z^2+j_{\mu,n}^2}.
\]
The recurrence
\[
I_\mu'(z)=I_{\mu-1}(z)-\frac{\mu}{z}I_\mu(z)
\]
then yields
\[
W_\nu(z)
=z\frac{I_{\mu-1}(z)}{I_\mu(z)}
=2\mu+2\sum_{n=1}^{\infty}\frac{z^2}{z^2+j_{\mu,n}^2}.
\]

Set \(s=z^2\). Then
\[
G_\nu(s):=W_\nu(\sqrt s)
=2(\nu+1)+2\sum_{n=1}^{\infty}\frac{s}{s+j_{\nu+1,n}^2}.
\]

Termwise differentiation is justified because \(j_{\nu+1,n}\sim \pi n\), so the differentiated sums converge locally uniformly on \((0,\infty)\). Thus
\[
G_\nu'(s)
=2\sum_{n=1}^{\infty}
\frac{j_{\nu+1,n}^2}{(s+j_{\nu+1,n}^2)^2}.
\]
For every \(k\ge0\),
\[
(-1)^k\frac{d^k}{ds^k}
\frac{j_{\nu+1,n}^2}{(s+j_{\nu+1,n}^2)^2}
=
(k+1)!\frac{j_{\nu+1,n}^2}{(s+j_{\nu+1,n}^2)^{k+2}}
\ge0.
\]
The locally uniformly convergent sum preserves these inequalities, hence \(G_\nu'\) is completely monotone and \(G_\nu\) is a Bernstein function.

If \(0<\tau\le1/2\), then \(x^{2\tau}\) is a Bernstein function. By Bernstein-function composition closure,
\[
x\mapsto W_\nu(x^\tau)=G_\nu(x^{2\tau})
\]
is a Bernstein function on \((0,\infty)\).
\end{proof}

The next result applies the preceding method to the fixed source problem \textbf{APP-0032}.
\begin{corollary}[APP-0032: From's all-\(L\) Mills-ratio bound family is explicit]
\label{res:from-mills-all-l-alternating-r-bound-family}
Let \(P_0=1\), \(P_1=t\), \(P_{n+1}=tP_n+nP_{n-1}\), and \(B_0=0\), \(B_1=1\), \(B_{n+1}=nB_{n-1}-tB_n\). With \(U_L(t)\) as in the all-\(L\) moment-ratio bound, the standard normal Mills ratio satisfies, for even \(L\ge0\), \(r(t)>[B_{L+1}-U_LB_L]/[P_{L+1}+U_LP_L]\), and for odd \(L\ge1\), \(r(t)<[U_LB_L-B_{L+1}]/[P_{L+1}+U_LP_L]\), for all \(t\ge0\).

This resolves the fixed application label \textbf{APP-0032}: Solve From's open all-\(L\) Mills-ratio bound problem from Remark 6.5.
\emph{Source reference.} \cite{app-app-0032-source}.
\end{corollary}
\begin{proof}
Set
\[
m_L(t)=\frac{M_{L+1}(t)}{M_L(t)}.
\]

The recurrence gives
\[
\frac{M_{L+2}}{M_L}=L+1-tm_L,
\]
\[
\frac{M_{L+3}}{M_L}=(t^2+L+2)m_L-t(L+1),
\]
and
\[
\frac{M_{L+4}}{M_L}=(L+1)(t^2+L+3)-t(t^2+2L+5)m_L.
\]

Substitution into the determinant inequality and expansion give
\[
(L+1)(t^2+4L+6)-t(t^2+4L+7)m_L-(t^2+4L+8)m_L^2>0.
\]

Equivalently,
\[
(t^2+4L+8)m_L^2+t(t^2+4L+7)m_L-(L+1)(t^2+4L+6)<0.
\]

Since \(m_L>0\), the positive root yields the uniform strict bound
\[
m_L(t)<U_L(t),
\]
where
\[
U_L(t)=
\frac{
\sqrt{t^2(t^2+4L+7)^2+4(L+1)(t^2+4L+8)(t^2+4L+6)}
-t(t^2+4L+7)
}
{2(t^2+4L+8)}.
\]

This is the determinant-to-bound extractor. It is a single theorem for every \(L\ge0\).

Define polynomials \(P_n(t)\) by
\[
P_0=1,\qquad P_1=t,\qquad P_{n+1}=tP_n+nP_{n-1}.
\]
Define \(B_n(t)\) by
\[
B_0=0,\qquad B_1=1,\qquad B_{n+1}=nB_{n-1}-tB_n.
\]

Then the same recurrence gives
\[
M_n(t)=(-1)^nP_n(t)r(t)+B_n(t).
\]

The ratio bound \(M_{L+1}<U_LM_L\) becomes
\[
\left((-1)^{L+1}P_{L+1}-(-1)^L U_LP_L\right)r
<U_LB_L-B_{L+1}.
\]

Since \(P_n(t)>0\) for \(t>0\) and \(P_{2j}(0)>0\), the coefficient has sign \((-1)^{L+1}\). Thus the bounds alternate:

For even \(L\ge0\),
\[
r(t)>
\frac{B_{L+1}(t)-U_L(t)B_L(t)}
{P_{L+1}(t)+U_L(t)P_L(t)}.
\]

For odd \(L\ge1\),
\[
r(t)<
\frac{U_L(t)B_L(t)-B_{L+1}(t)}
{P_{L+1}(t)+U_L(t)P_L(t)}.
\]

These formulas are valid at \(t=0\) by direct evaluation of the recurrences and for \(t>0\) by the positivity of \(M_L\).
\end{proof}

The next result applies the preceding method to the fixed source problem \textbf{APP-0016}.
\begin{corollary}[APP-0016: From's Mills-ratio bound extends to every derivative level]
\label{res:from-mills-explicit-bound-family-frontier}
After the Mills-ratio Laplace/complete-monotone normal form, derivative recurrence, and all-\(L\) determinant bridge, the remaining From frontier is solved by a closed, explicit alternating upper/lower bound family valid uniformly for all \(L\ge1\).

This resolves the fixed application label \textbf{APP-0016}: Lift the published Mills-ratio bound to the full derivative hierarchy.
\emph{Source reference.} \cite{app-app-0016-source}.
\end{corollary}
\begin{proof}
The density defining \(M_L\) is a positive moment density.  Its order-four
Hankel-minor consequence gives
\[
M_{L+4}M_L-4M_{L+3}M_{L+1}+3M_{L+2}^2>0.
\]
Using the recurrence \(M_{n+1}=nM_{n-1}-tM_n\), this determinant inequality
reduces to a quadratic inequality for \(m_L\); solving the admissible positive
root gives \(m_L<U_L\).  Iterating the same recurrence expresses \(M_L\) in the
form
\[
M_L=(-1)^L(P_L r-B_L),
\]
and substituting the level-\(L\) quotient bound gives the displayed alternating
upper and lower bounds for \(r\).
\end{proof}
